\begin{lemma}[Connectivity after finitely many full deletions]\label{lem:connectivity}
Let \(U\) be a plane topological disk and let \(C_1,\dots,C_N\Subset U\) be pairwise disjoint full continua.  Then
\[
 U\setminus\bigcup_{i=1}^N C_i
\]
is path connected.
\end{lemma}

\begin{proof}
Applying Lemma \ref{lem:jordan-basis} twice inside the disjoint neighborhoods used in Corollary \ref{cor:isolating-disks}, choose pairwise-disjoint Jordan domains
\[
 C_i\subset D_i^-,\qquad \overline{D_i^-}\subset D_i^+,\qquad \overline{D_i^+}\subset U.
\]
The finitely connected plane domain \(U_0=U\setminus\bigcup_{i=1}^N\overline{D_i^-}\) is connected: after a homeomorphism \(U\cong\mathbb D\), removing finitely many pairwise-disjoint closed Jordan disks compactly contained in \(\mathbb D\) leaves a disk with finitely many holes.  As an open connected subset of the plane, \(U_0\) is path connected.

For each \(i\), the open set \(D_i^+\setminus C_i\) is connected.  Indeed \(\mathbb C\setminus C_i\) is connected by fullness, and its intersection with the Jordan disk \(D_i^+\) cannot have a component separated from the outer collar of \(\partial D_i^+\): such a component would be a bounded complementary component of \(C_i\), contradicting fullness.  Hence \(D_i^+\setminus C_i\) is path connected.  It contains the nonempty annulus \(D_i^+\setminus\overline{D_i^-}\), which is also contained in \(U_0\); thus every \(D_i^+\setminus C_i\) has nonempty connected overlap with \(U_0\).  Their union is path connected and equals
\[
 U_0\cup\bigcup_{i=1}^N\bigl(D_i^+\setminus C_i\bigr)=U\setminus\bigcup_{i=1}^N C_i. \qedhere
\]
\end{proof}